\chapter*{How to use this book}
\addcontentsline{toc}{chapter}{How to use this book}

Each chapter moves from a concrete problem to concepts, formal reasoning, implementation, limitations, and exercises. Read the first explanation quickly, then slow down when the notation or code appears. A useful test of understanding is whether you can explain the idea without repeating the book's example.

The recurring case is not a single prescribed project. Treat it as a laboratory. Replace its public-service setting with a domain you know, but preserve the discipline of stating stakeholders, assumptions, data, interfaces, failure modes, and evidence.

\section*{A practical study loop}
\begin{enumerate}
\item Read the motivating problem and write down what is still unknown.
\item Run the smallest executable example before changing it.
\item Draw the system or model in your own notation.
\item Solve the conceptual exercises without looking at the solutions.
\item Modify the example to violate one assumption and observe the failure.
\item Record the decision, evidence, and limitation in a project log.
\end{enumerate}

The exercises are grouped as conceptual, technical, design, programming, research, and mini-project work. In collaborative work, distribute responsibility for implementation, modelling, user research, testing, and writing, but rotate roles so that everyone can explain the whole system and its limitations.

\warningbox{A passing test suite is evidence about the tested cases. It is not proof that the system is correct in every environment. The same distinction applies to user studies, simulations, and machine-learning evaluations.}
